% marc_aaai.tex — AAAI-2027 main-track submission (ported from the AAAI-2026 draft).
%
% Format: official AAAI-27 author kit (aaai2027.sty / aaai2027.bst, in this directory),
% wired in below via \usepackage[submission]{aaai2027}; the style sets
% \bibliographystyle{aaai2027} itself.
% Content rule: port + cut of marc.tex — every number identical, all \citep keys kept.
% The technical appendix (per-family tables, transfer, K-scaling) is marc_aaai_appendix.tex,
% input after the bibliography; it is intended as supplementary material.

\documentclass[letterpaper]{article} % DO NOT CHANGE THIS
\usepackage[submission]{aaai2027}  % DO NOT CHANGE THIS
% Fonts are loaded by aaai2027.sty (newtxtext, helvet, courier). The 2027 kit forbids
% adding times/helvet/courier or any other font package, so the 2026 preamble's font
% block and its XeTeX fontspec shim have been removed.
\usepackage[hyphens]{url}  % DO NOT CHANGE THIS
\usepackage{graphicx} % DO NOT CHANGE THIS
\urlstyle{rm} % DO NOT CHANGE THIS
\def\UrlFont{\rm}  % DO NOT CHANGE THIS
\usepackage{natbib}  % DO NOT CHANGE THIS AND DO NOT ADD ANY OPTIONS TO IT
\usepackage{caption} % DO NOT CHANGE THIS AND DO NOT ADD ANY OPTIONS TO IT
\frenchspacing  % DO NOT CHANGE THIS
\usepackage{booktabs}
\usepackage{amsmath,amssymb}
\pdfinfo{
/TemplateVersion (2027.1)
}

% Title and abstract track the OpenReview submission record (submission 49650); edits since
% then are confined to restoring the math the web form dropped and two accuracy fixes noted
% in the repo history.
\title{When Do Learned Diffusion Proposals Help Constraint Solving?\\
A Controlled Study on Continuous Algebraic Systems}
% Submission is double-blind: keep the anonymous line active until camera-ready.
% Camera-ready author block (equal contribution), swap in after acceptance:
%   \author{Quang Bui\equalcontrib, Sparsh Roy\equalcontrib,
%           Akash Gundimeda\equalcontrib, Davin Yin\equalcontrib}
%   \affiliations{SAID Laboratory}
\author{Anonymous submission}
\date{}

\begin{document}
\maketitle

\begin{abstract}
Evaluations of diffusion-based proposal models for constraint solving rarely include one key
control: random multi-start under the same refinement budget. When applied to our system,
this control sharply curtails our headline claim. MARC turns a continuous algebraic
constraint system into a factor graph, over which a graph-neural diffusion denoiser proposes
assignments, descent on an exact computer-algebra energy polishes them, and an exact symbolic
checker certifies solutions. Does the learned proposal improve on random multi-start at
choosing satisfying assignments? Only narrowly, in a predictable regime. Across trapped
low-dimensional families it ties with random restart, but dominates in high dimension, where
random search fails. Once variables couple, the advantage is gone. Since all methods share
the same polishing and checker, best-of-$K$ random multi-start succeeds with probability
exactly $1-(1-q(n))^K$, where $q(n)$ is single-start reachability; a single measured constant,
with no free parameters, reproduces the entire curve (mean absolute error 0.012). The narrow
favorable regime is not specific to our synthetic families: across eight real-world systems
in robotics, positioning, optimization, and algebra, classical multi-start solved all eight,
but none were in the learning-favorable regime. What classical solvers lack is discrete
choice: which structural augmentation renders an unsolvable system solvable. There, a
candidate-conditioned repair ranker outperforms random (0.997 vs 0.236 balanced nonlinear
menu accuracy; $p<10^{-70}$; $0.982\pm0.006$ across seeds), and beats a budget-matched
per-candidate probe on accuracy and cost. We delineate the regimes where learned proposals
improve solvers, and show that learning pays on the discrete decision, where it reaches the
exhaustive-search ceiling at a fraction of the calls.
\end{abstract}

\section{Introduction}
\label{sec:intro}

Solving a system of algebraic constraints requires a pair of decisions, one continuous and
one discrete. The former asks which values satisfy the constraints; it is differentiable and
has been studied for six decades in numerical analysis. The latter asks which representation
renders those values tractable to find: which auxiliary variables, substitutions, or defining
relations should be introduced. Classical numerical solvers take the representation as given,
making the structural choice only through enumeration; on the value decision they are the
method to beat. This paper evaluates what a learned component is worth in each, under one
controlled protocol. Learning has demonstrable value on the structural decision, where
enumeration is the classical alternative, and much less value on the value decision, where it
pays only inside a regime we can characterize precisely.

The evaluation is necessary because the common comparison between learned and classical
solvers conflates two contributions. A learned proposal plus refinement is better than a cold
start, which rewards the model for the value of diverse initialization. The control that
isolates that effect is random multi-start at equivalent polish and budget. Applied to our own
system it shrinks our headline claim, and finding where the learned component does pay is this
paper's subject.

MARC is an orchestrator: a diffusion denoiser proposes assignments over a factor graph, an
exact CAS supplies the residuals and gradients of a residual-based energy, deterministic
descent on that energy polishes each proposal, and a two-stage symbolic checker is the sole
acceptance criterion. The division of labor between a neural proposer and an exact symbolic
engine is the one AlphaGeometry demonstrated in olympiad geometry
\citep{trinh2024alphageometry}. The denoiser is a standard graph diffusion model and we
claim no novelty for it. MARC is the instrument of this study, not its contribution.

The result is a characterization with a single clear positive in it, obtained through a
sequence of progressively stronger controls; both are contributions. First, the learned
component is worth its cost on the structural decision. A repair ranker that conditions on
each candidate-augmented graph, rather than on the instance alone, clears its matched controls
on certificate-grade nonlinear menus and reaches the exhaustive-enumeration ceiling at a small
fraction of the calls. Second, that decision keeps its value away from benchmarks of our own
design: on hardened variants of named real systems, a construction derived from the givens and
selected on held-out failures repairs failures that no budget-matched restart control reaches.
Third, on the value decision the missing control shrinks the claim to a single regime, which
we then state as a law. A parameter-free factorization in the measured single-start
reachability $q(n)$ reproduces the separable family's best-of-$K$ curve and identifies the
conditions under which a learned proposal can win at all. Fourth, the same discipline applied
to our own pilot data uncovers a trap in how repair studies define their populations.
Admitting an instance because a stochastic solver failed on it once lets in instances that
succeed under fresh randomness, and it grades every arm on a stream correlated with the
criterion that selected them, so any intervention inherits apparent lift. We state the remedy
as a five-step protocol, two-stream failure selection, that applies to any study whose
population is defined by a stochastic solver's failures.

We also contribute the substrate itself, a pre-registered entrapment result, a partial
cross-family transfer result, and a MATH-benchmark scope measurement showing that
autoformalization, not solving, is the binding constraint. Every rate is reported with its
$N$ and a 95\% Wilson interval or $z$-test, and negative results are reported in full.

\section{Related Work}
\label{sec:related}

\paragraph{Multistart global optimization.}
Nothing in the classical half of our factorization law is new. Best-of-$K$ multistart succeeds
with probability $1-(1-q)^K$ in the single-start reachability $q$, and basin structure governs
the cost of stochastic global search \citep{rinnooykan1987clustering,
rinnooykan1987multilevel}; that annealed noise escapes minima where descent stalls is likewise
textbook \citep{welling2011sgld, bras2021langevin, regularized2025langevin}. Our entrapment
result is confirmation on this substrate, not novelty. What we add is the diagnostic: the
measured $\log q(n)$ slope, one constant reproducing the best-of-8 curve parameter-free, and
the two-condition dissection (collapse \emph{and} separability) that decides when a learned
proposal can beat the multistart bound.

\paragraph{Learned initialization and amortized optimization.}
A model mapping an instance to a solver starting point is a learned warm start:
\citet{amos2023amortized} surveys amortized optimization, learned starts accelerate AC
optimal power flow \citep{baker2019warmstart}, reinforcement learning tunes QP solvers
\citep{ichnowski2021rlqp}, and amortized Langevin inference \citep{taniguchi2022langevin} and
diffusion proposals with bootstrapped refinement \citep{bootstrapped2025refinement} are the
nearest neural instances. Our value-proposal negative applies directly to this line of work. Its
standard evaluation is learned start versus cold start. Under random multi-start at the
same refinement budget, a control this line does not run, our own advantage survives
only where solutions factorize per variable and dimension defeats random search.

\paragraph{Learning for combinatorial optimization.}
Two branches of neural CO border this work. One is graph diffusion solvers for NP-complete
problems: DIFUSCO \citep{sun2023difusco} and its unsupervised \citep{sanokowski2024diffuco},
trajectory \citep{li2024constraintaware}, and vehicle-routing \citep{anon2026vrpdiffusion}
successors. The other is learning-to-branch, from strong-branching surrogates \citep{khalil2016branch}
to bipartite-GNN branching \citep{gasse2019exact}, surveyed by \citet{bengio2021mlco}. MARC
differs in problem class and in where the verifier sits: continuous algebraic systems, an
exact CAS supplying residuals, gradients, and the sole acceptance gate, and a learned discrete
decision that is a pre-solve repair rather than in-tree branching. Neither branch reports the
budget-matched multistart control this study is built around.

\paragraph{Weak baselines as a known failure mode.}
That a learned method's advantage can be an artifact of its baseline is established in
neural combinatorial optimization. \citet{nath2024maxcutbenchmark} find published learned
heuristics for maximum cut that do not beat naive greedy, and that replacing one method's GNN
with linear regression leaves it intact; \citet{nath2023mightiest} report simple Tabu Search
beating learned local search. \citet{rodrigues2026budgetmatched} isolate the same pattern
outside combinatorial search. We claim neither observation nor remedy as new. What we add is that one polish and one
checker make Eq.~\eqref{eq:bestofk} the exact bound a learned proposal must clear, and the
measured $\log q(n)$ slope says in advance which families leave room above it.

\paragraph{Per-instance algorithm selection.}
The repair ranker is, at bottom, per-instance selection: given an
instance and $K$ discrete options, pick the one under which a downstream solver succeeds.
The problem is fifty years old \citep{rice1976algorithm}, was made standard practice by
SATzilla \citep{xu2008satzilla}, and is surveyed in \citet{kerschke2019survey}; AlphaGeometry
\citep{trinh2024alphageometry} is the precedent for learning
the structural decision itself, and discrete diffusion \citep{austin2021d3pm,
vignac2023digress} supplied the formalism for our withdrawn predecessor policy. Three things
here are not in that line: certificate-grade menu semantics (``exactly
one solvable option'' is a CAS theorem, not a budget-relative claim), candidate conditioning
(the ranker encodes each candidate-augmented graph, not instance features alone), and
controls that price the classical recourse, namely budget-matched restarts and a
per-candidate solver probe, rather than a single default algorithm.

\paragraph{Distance geometry.}
The geometric family below, chains of points pinned by squared-distance
constraints, is a distance geometry problem with known discrete structure: each new point
lies on an intersection of circles, the per-point reflection ambiguity is the discrete branch
that branch-and-prune enumerates, and the solution set carries a group of partial reflections
\citep{lavor2012dmdgp, liberti2014distance}. We claim none of that structure. Those branches
are exactly the ambiguities our reachability measurement counts. What we add is measurement under shared machinery: the reachability slope under one fixed
polish, and a trained proposal that ties restart-matched random search there anyway.

\section{Method}
\label{sec:method}

A problem instance is a factor graph $G = (V, F)$ over variable nodes $V = \{x_1, \dots,
x_n\}$ and factor nodes $F = \{g_1, \dots, g_m\}$ carrying symbolic expressions: a bare $g$
denotes the equality $g = 0$, relational expressions denote non-strict inequalities. The
residual $r_j(x)$ is $|g_j(x)|$ for equalities and the hinge $\max(0, g_j(x))$ for
inequalities, so $r_j(x) = 0$ exactly when constraint $j$ holds. The energy is
\begin{equation}
E(x) \;=\; \tfrac{1}{2} \sum_{j=1}^{m} r_j(x)^2 .
\end{equation}
Residuals, energy, and $\nabla E$ are computed exactly by a CAS (SymPy) and compiled once per
graph. Acceptance is a two-stage gate: a numeric stage rejects fast if any factor's violation
exceeds a tolerance, and a symbolic stage snaps the candidate to nearby exact rationals and
re-checks every constraint exactly. On the synthetic families the symbolic stage is
authoritative. On the real systems, whose roots are irrational, no exact rational exists to snap to
and the numeric stage stands alone. This gate is the only acceptance criterion in every
experiment and the only reward source in training.

\paragraph{The learned proposal.}
The forward process is standard Gaussian diffusion over variable values with a cosine noise
schedule over $T = 1000$ steps. The denoiser is a bipartite message-passing GNN: variable
nodes are encoded from their noisy value and type, factor nodes from their type,
residual, and a sinusoidal timestep embedding; $L$ rounds of message
passing are followed by a per-variable head predicting the noise
$\hat{\varepsilon}$. Sampling runs DDIM reverse steps with energy-gradient guidance,
\begin{equation}
\hat{s} \;=\; s_\theta(x_t, t, G) \;-\; \lambda_t \nabla E(x_t),
\end{equation}
where $s_\theta$ is the score implied by the predicted noise ($s_\theta \propto
-\hat{\varepsilon}$); the exact CAS gradient steers the reverse process toward feasibility.
One architectural detail proved necessary: a linear skip from incident-factor constants
to the variable's output, without which per-round LayerNorm washes out pinned values
(mean absolute root error 5.4 versus 0.9; Appendix A of the supplement).

\paragraph{Propose and polish, and the baseline battery.}
The system solver is a hybrid: the diffusion model proposes $K$ candidate assignments
(best-of-8 throughout); each is polished by deterministic descent on $E$; the first candidate
to pass the checker is the answer. Each baseline isolates one ingredient: deterministic energy descent ($\sigma_k \equiv 0$ in
$x^{(k+1)} = x^{(k)} - \eta \nabla E(x^{(k)}) + \sigma_k \xi_k$) stalls where
$\nabla E = 0$ but $E > 0$; annealed Langevin ($\sigma_k$ decreasing) is the classical
stochastic solver; the mean-prior control tests whether the model does more than reproduce
an average solution. The decisive control is random multi-start: $K$ initializations drawn
uniformly from a family-matched box, polished identically, same budget, no learning.

\section{Experiments}
\label{sec:experiments}

\subsection{Protocol}
\label{sec:entrapment}

All problem families are procedurally generated with held-out test instances. Unless stated
otherwise, solve rates are best-of-8 over $N = 60$ held-out instances per family, intervals
are 95\% Wilson intervals, and comparisons are one-sided two-proportion $z$-tests of
learned $>$ baseline. Before any comparison we
verified that the learned solver converges at all: after five implementation faults were found
and fixed (Appendix A of the supplement), the solve rate on convex linear systems
went from 0 to 1.000 in-distribution and held out. Convex families are saturated by every
method and carry no comparative signal. All results below are non-convex.

A pre-registered check confirmed that stochasticity matters in
this search space. Deterministic energy descent is trapped on
every one of 200 non-convex instances (rate 1.000), and annealed Langevin cuts entrapment to
0.475, a reduction of $0.525 \pm 0.109$ whose 95\% interval excludes zero (5 seeds). This is
textbook behavior, and we present it as a confirmed premise rather than a finding.
% provenance: RESULTS.md R2 (entrapment table folded into prose; full table was tab:entrapment)

\paragraph{Two-stream failure selection.}
One protocol element is stated in reusable form because a later result lives or dies by it.
Repair-style studies build their population from instances a stochastic pipeline fails.
Conditioning on one bad draw admits instances that succeed under fresh randomness and grades
every control on a stream correlated with the failure criterion, so any intervention inherits
apparent lift. The discipline: (i) admit an
instance only if the direct solve fails on two independent streams; (ii) label interventions
on a third, majority-voted when the labels train a model; (iii) grade all arms on a fourth,
common across arms; (iv) field a matched-budget restart control beside the interventions;
(v) before crediting a learnable selection target, screen each candidate at the target budget
on fresh streams and grade the argmax on a held-out one. If that per-instance measurement
cannot beat the matched-budget control, the advantage is portfolio breadth, not selection a
model could amortize. On our geometry population single-stream selection put the construction ceiling at
$0.63$--$0.74$ against $0.09$--$0.27$ for matched restarts, and two-stream selection with a
common grading stream collapsed the margin, most of it the control returning to its true
rate.

\subsection{Diverse starts and polish account for the hybrid's performance}
\label{sec:hybrid}

On four non-convex families where deterministic descent is fully trapped (best-of-8, 60
held-out instances per family), cold-start refinement solves nothing and Langevin noise helps
only somewhat, while both the random-init control and the learned hybrid do far better (full
battery in supplementary Table S1). What matters is the comparison of learned
against random, and there the learned proposal records no win on any of the four families:
three ties and one complete failure on CircleLine (0.000 where random reaches 0.200).
Multi-start Levenberg--Marquardt saturates all four, so the informative contrast is learned
against random, not learned against cold-start Langevin.
% provenance: RESULTS.md R3
Without the random control the ablation reads as ``the
learned solver beats classical refinement''; with it, random multi-start plus deterministic
polish beats cold-start Langevin and the proposal adds nothing, because the solutions are
small integers that random restart hits readily and there is nothing to
amortize. What carries this result is the restart-and-polish recipe, not the denoiser. Cross-family transfer is likewise
partial. The cross-trained hybrid reaches 0.683 on two of four held-out families
($p < 10^{-4}$) and solves nothing on the other two
(Appendix A of the supplement).

\subsection{Dimension scaling: where the learned proposal wins}
\label{sec:scaling}

Those families give random restart an easy target. To test the amortization argument we
use bundled non-convex traps whose roots vary per instance over a wide signed range
($\pm[3,8]$), so neither a fixed prior nor a lucky draw suffices (best-of-8, unified-v2
protocol: one shared polish and one checker for every arm, $N{=}40$ per cell; the full battery
is supplementary Table S3, and Figure~\ref{fig:law} carries the two arms that
matter).
% provenance: RESULTS.md R15 (methodology unified-v2, results/p_scaling/scaling.json)
There is a clean crossover by $n = 3$. At $n = 1$ random restart wins, 1.000 against 0.950.
Its solution space is small enough to brute-force, so learning buys nothing. From $n = 3$
random restart collapses (0.075, then 0.000), because a random initialization must land in
all $n$ basins at once. The learned
proposal holds essentially flat where random fails: 0.950 at $n = 2$, 0.975 at $n = 3$, 0.925
at $n = 4$. That is the ceiling the law below takes as its input.
Langevin and the mean-prior sit at or near zero from $n = 3$. This is the amortized-inference
result, a per-instance proposal replacing search whose cost is exponential in dimension,
with one caveat: the learned arm is seed-unstable at intermediate $n$ (single seed
0.250 at $n = 6$; three seeds read $0.983 \pm 0.014$ there but $0.658 \pm 0.484$ at
$n = 4$). The law's larger-sample measurement below reads random lower at low $n$
and places the crossover earlier, so we report the geometric collapse, not a
single $n^*$.

Two objections apply here: that this is one designed family, and that random multi-start is a
weak baseline. Rerun on three structurally different separable families with a
Levenberg--Marquardt arm (analytic Jacobian, eight Gaussian
multistarts), the crossover replicates, and LM is not a way out, since it must also hit all
$n$ independent basins and collapses along the same $v^n$ curve.
% provenance: RESULTS.md R27 + PROVENANCE R27 (crossover_families.json)
On two of the three the learned
proposal significantly beats \emph{both} random restart and LM at high dimension (learned
$1.000$ against $0.000$ for both at $n = 6$ on the baseline family). Crossover is
therefore a property of separable structure, not of one construction.
On the third (two spurious wells) the small denoiser did not learn the harder
marginals and collapsed with the classical methods. That sharpens the boundary. Learning
wins only where it can actually amortize per-variable marginals.

\subsection{A factorization law predicts both results}
\label{sec:law}

The crossover above and the coupled negative reported next are two outcomes of one mechanism,
statable as a law and testable without fitting anything. Every method in
this study shares one polish operator and one acceptance checker and differs only in where
the polish starts. Define the single-start reachability
% provenance: RESULTS.md R9 + paper/notes/crossover_law.md §1-2
\begin{equation}
q(n) \;=\; \Pr_{x_0 \sim U}\big[\,\text{polish}(x_0)\ \text{is accepted}\,\big],
\end{equation}
a geometric property of the problem family and the polish operator, with no
learning in it. Random multi-start with budget $K$ succeeds iff at least one of $K$ i.i.d.\
starts is accepted, so
\begin{equation}
P_{\text{random}}(n;K) \;=\; 1 - \big(1 - q(n)\big)^{K} .
\label{eq:bestofk}
\end{equation}
The identity is exact for a fixed instance; across a family we use the instance-averaged
version, and we bound the Jensen gap this opens by validating against a self-measured
best-of-8 curve rather than assuming it negligible.
It already implies that a hybrid beating cold-start Langevin has shown nothing. It must beat
Eq.~\eqref{eq:bestofk} at the same $K$.

How $q(n)$ scales with dimension is decided by whether the acceptance basins factorize across
variables. If each factor touches one variable, the polish is coordinate-decoupled and a start is
accepted only when every coordinate independently lands in its root basin, so $q(n) = v^n$
with $v := q(1)$: $\log q$ is linear in $n$ and random search needs $\Theta(v^{-n})$ starts
to hold a fixed success rate. If factors couple, the solution is a joint object and $q(n)$
need not decay geometrically at all. We measured $q(n)$ directly by single-start polish on 600 fresh
instances per $n$ (Wilson intervals), using the same generators, polish, and checker as the
solve-rate experiments. supplementary Table S4 tabulates the dichotomy.

% provenance: paper/tex/figures/fig_crossover_theory.pdf, RESULTS.md R9
\begin{figure*}[t]
\centering
\includegraphics[width=\textwidth]{figures/fig_crossover_theory.pdf}
\caption{The factorization law, measured. Left: $\log q(n)$ against $n$. The separable
family is a line (slope $-1.03$, $R^2 = 0.98$), the coupled bilinear family is nearly flat
($-0.13$), and the geometry family, though syntactically coupled, collapses ($-0.77$,
$R^2 = 0.999$). Right: the best-of-8 random-restart curve predicted parameter-free from
$v = 0.27$ via Eq.~\eqref{eq:bestofk} against the measured curve (MAE 0.012), with the
learned hybrid overlaid; the dotted line is the crossover dimension predicted from $v$.}
\label{fig:law}
\end{figure*}

The separable family behaves exactly as the law requires: $\log q(n)$ is linear with slope
$-1.03$ ($R^2 = 0.98$). Substituting the directly measured $v = q(1) = 0.27$ into
Eq.~\eqref{eq:bestofk} reproduces the entire best-of-8 random-restart curve with no free
parameters. It predicts 0.919,
0.454, 0.147, 0.042, 0.003 against measured 0.902, 0.467, 0.162, 0.030, 0.002 at
$n = 1, 2, 3, 4, 6$, a mean absolute error of 0.012
% provenance: RESULTS.md R9 + crossover_law.md §5 (crossover_theory.json)
(Figure~\ref{fig:law}; 600 fresh instances, a larger and separately drawn sample than
supplementary Table S3's $N = 40$ random arm). A budget-fair
reading is the expected number of restarts $1/q(n)$: 3.7 to 600 over the separable
dimensions against 2.0--4.3 on the coupled family. Any fixed budget is exhausted on the
former, while on the latter random search never collapses and a learned proposal has
nothing to amortize. An oracle control
closes the mechanism: sampling each variable from its \emph{true} per-variable marginal
(pooled from disjoint instances) and polishing identically also ties random restart at every
$n \ge 3$ on the coupled family (0/4 wins). Sampling coordinates independently, it
bounds product-of-marginals proposals only. Population-marginal information is exhausted
there, and our denoiser behaves like a marginal sampler rather than an instance-conditional
joint one.

Separability is sufficient for geometric collapse but not necessary, so the diagnostic is the
measured slope, not the syntactic label. A real-valued geometric domain makes this concrete:
chains of unknown points ($n = 2k$ coordinates) with squared-distance constraints to fixed
anchors and between consecutive points, a non-convex quartic energy whose solutions the
checker accepts exactly. The family is syntactically coupled, yet its reachability collapses:
slope $-0.77$ ($R^2 = 0.999$), $q(n) = 0.653$, 0.147, 0.027, 0.007 at $n = 2, 4, 6, 8$
% provenance: RESULTS.md R9 (geometry) + crossover_law.md §5a
(Figure~\ref{fig:law}, left). Each point's two-circle subproblem has a reflection
ambiguity, so a random start rarely places every point in the right basin. Collapse alone
would suggest an amortized proposal wins here. Training a denoiser on this family under the
identical protocol does not bear that out. The learned proposal ties random restart at every
chain length and collapses with it (both 0.625, 0.175, 0.025, 0.000 at $n = 2, 4, 6, 8$,
$N = 40$ per length; zero of four significant wins).
% provenance: RESULTS.md R9 (geometry) + PROVENANCE R25 (pointchain_learned.json)
That negative is the informative result, because it separates two conditions the independent
family had conflated. Reachability collapse (condition one, read off the slope) determines
whether random search fails. It does here. Separability (condition two) determines whether the
denoiser has per-variable marginals to amortize. Geometry fails it, because each point's
coordinates are pinned only jointly. Both conditions are necessary, so a learned proposal
beats classical search only where reachability collapses \emph{and} the solution is
per-variable separable; given those, it wins when the denoiser fits the marginals.
Figure~\ref{fig:regime} places every measured family on the two axes of the law.

% provenance: paper/tex/figures/fig_regime_map.pdf (scripts/plot_regime_map.py);
% slopes from crossover_theory.json (R9), R27 slopes law-inverted from the LM arm of
% crossover_families.json; outcomes from scaling.json (R15), crossover_families.json (R27),
% coupled.json (R7), pointchain_learned.json (R25), real_systems.json (R26)
\begin{figure}[ht]
\centering
\includegraphics[width=\columnwidth]{figures/fig_regime_map.pdf}
\caption{The regime map. Each measured family sits at its measured $\log q(n)$ slope
(abscissa) in its solution-structure band (ordinate, categorical), colored by the measured
learned-vs-random outcome; slope provenance is in Appendix C of the supplement. The cells give the law
directly: learning wins only where reachability collapses \emph{and} solutions are
per-variable separable. The chained bilinear family fails collapse and ties; the geometry
point chains collapse but are coupled and tie, which was the falsifiable cell and it held. The
double-well tie inside the win cell is a capacity failure, not a regime failure; the
real-systems suite has no learned arm and no measured slope (nominal abscissa).}
\label{fig:regime}
\end{figure}

\subsection{Coupling removes the advantage}
\label{sec:coupled}

The law predicts where the crossover must disappear: when acceptance basins do not factorize,
$q(n)$ stays nearly flat, random restart never collapses, and a learned proposal has nothing
to amortize. The test is the coupled chained-bilinear family, $x_i + x_{i+1} = s_i$ and
$x_i \cdot x_{i+1} = p_i$, whose solution is a joint object and whose measured reachability
slope is $-0.13$ (supplementary Table S4), so what follows is a prediction checked, not
an unexplained negative. Table~\ref{tab:coupled} reports best-of-8 with 60 test instances per
$n$.

% provenance: RESULTS.md R7
\begin{table}[t]
\centering
\caption{Coupled chained-bilinear systems (best-of-8, 60 test per $n$). The learned proposal
records zero significant wins over random restart in five comparisons.}
\label{tab:coupled}
\footnotesize
\setlength{\tabcolsep}{4pt}
\begin{tabular}{ccccl}
\toprule
$n$ & Langevin & random & learned & learned $>$ random? \\
\midrule
2 & 0.167 & 0.483 & 0.233 & no (loses) \\
3 & 0.100 & 0.600 & 0.533 & no \\
4 & 0.050 & 0.517 & 0.517 & no (tie) \\
6 & 0.000 & 0.367 & 0.333 & no \\
8 & 0.000 & 0.467 & 0.483 & no \\
\bottomrule
\end{tabular}
\end{table}

Across all five dimensions the learned proposal ties or loses, with 0 significant wins. That high-dimensional advantage
was an independence artifact. It needed per-variable-separable
solutions the model can memorize as marginals, and a random restart forced to hit all $n$
basins by chance. Coupling removes both conditions, and the proposal adds nothing over
random search plus refinement. This bounds the claim.

\paragraph{Standard real systems.}
Every family above is procedurally generated, so we also ran the solver battery on eight
recognized test problems encoded once as factor graphs, graded by the numeric stage of the
checker ($\max_j |r_j(x)| < 10^{-6}$) since their roots are irrational
(Table~\ref{tab:real}).
% provenance: RESULTS.md R26 + real_systems.md (results/p_real/real_systems.json)
Multi-start Levenberg--Marquardt solves all eight; annealed Langevin solves one and random
restart with the gradient polish solves four. Where the gradient polish fails, always at single-start reachability $0.00$, the stronger
\emph{classical} polish fixes every case: the bottleneck is the finisher, not the proposal.
No system falls in the amortization regime, since all eight are low-dimensional and coupled,
so we do not run the learned component per system.

% provenance: RESULTS.md R26; results/p_real/real_systems.json (rows[].q_single_start)
\begin{table}[t]
\centering
\caption{The eight standard systems (best-of-8; numeric acceptance at $10^{-6}$).
$q$ is single-start reachability under the gradient polish, 200 draws; $\nabla$ marks the
stationary-point systems. Deterministic descent solves none. Every system the gradient polish misses sits at $q = 0.00$ and is
solved anyway by multi-start Levenberg--Marquardt: the bottleneck is the finisher, not the
proposal.}
\label{tab:real}
\footnotesize
\setlength{\tabcolsep}{3pt}
\begin{tabular}{lccccc}
\toprule
System & $n$ & $q$ & Langevin & random & LM \\
\midrule
Circle intersection    & 2 & 1.00 &            & \checkmark & \checkmark \\
Conic--line            & 2 & 1.00 &            & \checkmark & \checkmark \\
GPS trilateration      & 2 & 0.00 &            &            & \checkmark \\
Rosenbrock $\nabla$      & 2 & 0.00 &          &            & \checkmark \\
Himmelblau $\nabla$      & 2 & 0.00 &          &            & \checkmark \\
2R inverse kinematics  & 4 & 0.98 & \checkmark & \checkmark & \checkmark \\
3R inverse kinematics  & 6 & 0.00 &            &            & \checkmark \\
Cyclic-4               & 4 & 0.38 &            & \checkmark & \checkmark \\
\bottomrule
\end{tabular}
\end{table}

\subsection{Relocating the learned component: structural repair beats its controls}
\label{sec:repair}

Coupling closes the route where the network predicts values, and the real-systems suite
confirms classical search owns that game. The one decision classical solvers cannot make
is discrete and prior to any solve: \emph{which} structure-changing augmentation turns an
unsolvable graph into a solvable one. We move learning there. A candidate-conditioned
ranker applies each of $K$ proposed augmentations (here $K$ is the menu size, not the restart
budget of Eq.~\eqref{eq:bestofk}), encodes the resulting polynomial graph with operator-aware
node and edge features (degree, linear/square/cross participation, constants), and scores the
augmented graphs listwise; the highest-scoring repair receives one classical solve and the
exact checker remains the acceptance gate. A matched candidate-only control sees the same augmentation
recipe but no problem graph, and \texttt{random} chooses uniformly among the same $K$
candidates. The predecessor policy classified menu slots over one fixed graph encoding and fell to
chance on an unseen pattern. Candidate conditioning is the change that matters.
An ablation locates the signal: masking operator-identity features and
retraining leaves the ranker intact ($0.978$ $[0.957, 0.989]$ against $0.997$ unmasked,
$N{=}360$; reruns at fixed seed move the third decimal, so we read the interval). The signal lives in constants, magnitudes, and incidence read
\emph{jointly with the candidate}, not in operator flags.

% provenance: RESULTS.md R10 + PROVENANCE R20--R24; results/p_repair (Data Version 8)
\begin{table}[t]
\centering
\caption{Candidate-conditioned repair ranker against a candidate-only control and random
selection, on three generalization tests (Data Version 8; significance by exact paired McNemar in
the text). Nonlinear menus
carry exact CAS no-real-roots certificates for their distractors; every arm is graded by the
same reference solver that certified the data.}
\label{tab:repair}
\footnotesize
\setlength{\tabcolsep}{3.5pt}
\begin{tabular}{lrccc}
\toprule
Test & $N$ & ranker & cand.-only & random \\
\midrule
Balanced nonlinear      & 360 & \textbf{0.997} & 0.333 & 0.236 \\
Vieta $\to$ unseen rel.\ & 150 & \textbf{0.420} & 0.120 & 0.253 \\
Linear, unseen pattern  & 400 & \textbf{0.380} & 0.195 & 0.287 \\
\bottomrule
\end{tabular}
\end{table}

Table~\ref{tab:repair} reports the three generalization tests under a deliberately strict
protocol: one reference solver certifies the data, grades every arm,
and runs the end-to-end solves; nonlinear ``exactly one solvable option'' is an exact CAS theorem (a distractor proven to
have no real solution is unsolvable at any budget); gold and distractor parameters share one
support and one prior, so surviving surface-form signal is capped at the candidate-only
ceiling and anything above it must come from reading the problem graph. On the balanced nonlinear test that ceiling is $0.333$, above the $0.25$ chance floor and
far below the ranker, where exact paired McNemar gives $p = 3.3\times10^{-83}$ against
random (274 ranker-only correct versus 0) and $p = 1.1\times10^{-72}$ against the
candidate-only control. Transfer to an unseen nonlinear relation is partial but bidirectional ($0.420$ against
$0.253$ for vieta-trained on quad\_link, $0.393$ against $0.180$ reversed; $N{=}150$ each). On an unseen linear pattern the ranker reaches $0.380$ against $0.287$ ($N{=}400$), and
holding out each of the other two patterns gives $0.407$ and $0.450$ against $\sim$$0.23$
random, so the effect is not specific to one holdout. A single model
tested across all three at once reads $0.339$ against $0.249$ ($N{=}1{,}200$, $p =
7.8\times10^{-7}$). The linear signal is reported at its post-audit size: closing each discovered shortcut
lowered it ($0.565 \to 0.445 \to 0.380$ across data versions) while pushing the
candidate-only control to chance, so what remains is problem-graph reading, modest where
every operator is linear and decisive on certificate-grade nonlinear menus. Both
headline rows were retrained across three optimization seeds (11, 29, 47) with independent
per-seed evaluation draws. Nonlinear comes through optimization-robust ($0.982 \pm 0.006$).
The linear edge is not seed-stable ($0.317 \pm 0.069$, one seed at $0.227$ below the random
arm's $0.248 \pm 0.002$), consistent with reading the linear rows as mechanism evidence
rather than a deployable margin.

Learning is also cheap here where value diffusion was not: one forward pass per
candidate rather than a reverse-diffusion rollout. End to end, the nonlinear ranker's single
call solves $0.933$,
exactly the oracle and enumeration ceiling, where blind enumeration averages $2.62$ calls
($N{=}60$); linear $K{=}4$ solves $0.300$ against $0.227$ random with a $1.000$ ceiling
($N{=}300$). A cheap-probe control bounds the claim from the other side: a short-budget call on
every candidate solves at most $0.881$ of nonlinear menus at $4.7$ calls per instance, where
the ranker's single call solves $0.939$ ($N{=}360$). Learning beats probing on accuracy
and cost at once. Provably rootless distractors are unsolvable at any budget, and short
probes miss the gold. On linear menus the probe saturates and enumeration is
already perfect at $2.5$ calls. The $K = 4$ checkpoint transfers to larger menus without
retraining, but its advantage is gone by $K = 16$, and direct $K = 16$ training sits at chance (both negatives kept in
Appendix E of the supplement). ``Exactly one solvable option'' is exact for all linear menus (rank) and $99\%$ of nonlinear
test menus (CAS real-root nonexistence); the remainder carry a disclosed budget-relative
probe certificate.

% provenance: RESULTS.md R30; results/p_real_repair/real_repair{,_multiseed}.json
\paragraph{The structural decision on systems we did not design.}
Those menus are ours. We therefore hardened four of the named real-system classes into
parameterized variants and asked whether \emph{derived} constructions, functions of the
givens alone such as law-of-cosines lifts and signed Cayley--Menger cross-product pins,
repair the failures the classical reference cannot (supplementary Table S5 reports
all four classes). Two produce a two-stream failure population: far-side GPS trilateration ($0.848 \pm 0.020$
over three seed bases, $N{=}509$ failures) and a ghost-root conic--line intersection
($0.263 \pm 0.015$, $N{=}158$). On 3R inverse kinematics and far circles the classical solver
never fails, which we report as a negative rather than averaging in as a zero. On the two classes that do produce failures, one construction chosen on a disjoint half of
each pool repairs
\emph{every} held-out instance ($1.000 \pm 0.000$ across seeds; pooled Wilson $[0.99, 1.00]$
and $[0.98, 1.00]$) against $0.433 \pm 0.049$ and $0.114 \pm 0.011$ for a restart control
matched to the \emph{full enumeration budget} over the vocabulary (McNemar
$p < 10^{-13}$; the control, not the treatment, carries the sampling variance). The same
construction wins in all six folds, which rules out its being the luckiest of $V$, and
saturation is
mechanical rather than fortunate: the selected pins delete the mirror basin and the ghost
root by construction. No learned model runs here. What transfers to systems we did not
design is the value of the structural decision, not the ranker.


\paragraph{Where the construction result does not extend.}
This result holds where the failure is a systematic attractor that one construction deletes
outright. Where failures are stochastic it does not, and the boundary is worth measuring.
On the pruned point chains, the discrete-branch setting of distance geometry, the
population definition decides the answer. Under two-stream selection ($N{=}367$ failures, three optimization seeds) the enumeration
ceiling is $0.692$ at $72.7$ restarts per instance, which plain restart scaling matches at
$+32$ ($0.725$). Majority-vote labels made the target learnable, and the ranker separates from random
($0.246 \pm 0.016$ versus $0.185 \pm 0.019$; Holm $p{=}1.3{\times}10^{-4}$) where
single-stream labels never did. That is not an edge over the controls: it ties the best
fixed construction ($0.259$) and loses to matched-budget restarts ($0.270$). The signal is real but
reflects the population prior; the problem graph contributes nothing. The probe initially
appeared promising: one restart per candidate solves $0.698$, matching the ceiling. But it spends the whole menu's budget and its accepts are diffuse, so it is a portfolio
sweep, not a choice of construction. A per-instance screen held to the reference budget
scores $0.199$ $[0.161, 0.243]$, below the restart control: ${\sim}50$ restarts of
measurement select worse than the prior. The learned arms landing on the prior is therefore expected. Scoring that same screen on
the streams that chose it reads $0.762$, which prices the selection-on-noise effect. A
construction pays when it removes a named failure mode outright; where none does,
constructions help only as a portfolio and ``which construction'' carries no signal a
selector can exploit.

\section{Limitations}
\label{sec:limitations}

Our positive result has a stated scope. The external anchor uses \emph{derived} constructions
and a held-out choice among them, not the learned ranker; its populations are hardened
variants chosen to expose a failure mode, and two of the four produce no failure population. We measure the ranker itself on menu-based repair over synthetic factor
graphs.

On the value side, the learned proposal's advantage requires per-variable-separable
solutions and a regime where random restart collapses. It gains nothing at $n = 1$ (random restart at ceiling), and single cells are seed-noisy, so
the useful window is the crossover region, not high dimension per se. Those families are also
block-decomposable: a solver that read off separability would scale linearly there too, so
the crossover is demonstrated against joint-start search, not block decomposition. Eq.~\eqref{eq:bestofk} predicts the random-restart curve, not the learned ceiling, which
is measured.

Several failures are specific. CircleLine is never solved by the learned proposal, transfer
succeeds on two of four held-out families, and adding CircleLine to the training mix
collapsed an otherwise recoverable transfer (0.70 to 0.00). The predecessor structure policy never
separated from random selection and its numbers remain withdrawn (evaluation seeds
overlapped checkpoint-selection validation seeds), as do the earlier v6/v7 ranker numbers
(a pin-prior leak and probe-artifact certificates, found and fixed). Only the nonlinear headline ($0.982 \pm 0.006$) is optimization-robust.
Finally, scope against real problem text is quantified: on a 48-problem sample of
MATH-500 \citep{hendrycks2021math} the template formalizer covers none of them
(breakdown in the supplement), so the binding constraint is autoformalization, not the
solver.
% provenance: RESULTS.md R6 (MATH prose); withdrawal provenance: RESULTS.md R8/R10

\section{Conclusion}
\label{sec:conclusion}

The decision classical solvers make only by enumeration is which structure to add, not which
values to try, and that is where the learned component pays: the repair ranker matches the
enumeration ceiling at a fraction of the calls where ``exactly one solvable option'' is a
theorem, and on hardened real systems a derived construction repairs failures the full
restart budget misses. On the value decision the protocol returns a characterization
instead. The learned proposal pays only where search is separable and dimension defeats
random restart, and the advantage disappears under coupling. The generalizable lesson is
about evaluation: populations defined by stochastic failure are artifacts of the draw unless
the protocol stabilizes them, at one extra solve per instance.



% AAAI-27 mandates this order: main content, content appendices, ethical statement,
% acknowledgments, references, supplementary material. Unnumbered statements go here,
% before the bibliography.
\section*{Use of AI Assistants}
A large language model assisted with the readability of this paper. All experiments,
analysis, and claims are the authors' own; no result, table, or figure was model-generated.

% aaai2027.sty already sets \bibliographystyle{aaai2027}; a second one breaks bibtex
\bibliography{refs}


% The technical appendix ships as a SEPARATE supplementary document,
% paper/tex/supplement.tex. AAAI-27 counts appendices bound into the main PDF
% against the 7 content pages and requires supplementary material to be
% uploaded separately, so it is no longer \input here.
\end{document}
